% quantitative_results.tex — standalone fragment, \input'd into main paper
% No preamble; assumes booktabs, amsmath, graphicx, hyperref loaded

\section{Quantitative Results}
\label{sec:quantitative_results}

This section presents the model's quantitative predictions. Section~\ref{sub:timing}
examines steady-state first-birth timing. Section~\ref{sub:scarcity} turns to persistent
housing scarcity. Section~\ref{sub:voter} compares voter preferences in the fertility and
NIMBY benchmarks. Section~\ref{sub:evidence} closes with brief supporting evidence.

%% ============================================================
\subsection{Steady-state first-birth timing}
\label{sub:timing}

The cleanest result from the fertility extension is a steady-state comparison across
house-price levels. Because parity is an explicit household state, the model can trace
first births directly rather than only tracking total or completed fertility.
Figure~\ref{fig:first_birth_timing} plots the three timing objects---mean age at first
birth, share of first births at age 30 or later, and the first-birth rate---against the
steady-state house price~$q$.

The central result is a comparison of fertility outcomes at the two political-equilibrium
prices. The NIMBY benchmark (fertility shut off) crosses at $q = 2.326$; the refreshed
fertility benchmark now crosses locally at $q = 1.750$. Table~\ref{tab:equilibrium_comparison}
reports the demographic consequences of this gap, evaluating the fertility model at those
two benchmark prices. Moving from the fertility-model crossing to the higher NIMBY
crossing price delays mean first birth by about 1.4~years, doubles the share of first
births at age 30 or later, raises the childless share by about 19.8 percentage points,
and lowers the first-birth rate from about 0.213 to about 0.163. The large completed-
fertility adjustment comes mainly through a collapse in the share of households with three
or more children, from about 41.7 to about 15.3~percent. These are still meaningful
differences in completed family size, but the cleaner and better-disciplined model object
is the timing margin itself: expensive housing pushes entry into parenthood later, and
that delay carries through into fewer realized births by age 50.

\begin{table}[htbp]
\centering
\caption{Fertility-model outcomes at the two benchmark crossing prices}
\label{tab:equilibrium_comparison}
\small
\begin{tabular}{l cc c}
\toprule
 & Fertility crossing & NIMBY crossing & \\
 & ($q = 1.750$) & ($q = 2.326$) & Difference \\
\midrule
\multicolumn{4}{l}{\textit{First-birth timing}} \\
Mean age at first birth & 26.22 & 27.60 & $+1.38$ years \\
Share of first births at 30+ & 0.202 & 0.405 & $+0.204$ \\
First-birth rate & 0.213 & 0.163 & $-0.050$ \\
\addlinespace
\multicolumn{4}{l}{\textit{Parity distribution at age 50}} \\
Childless & 0.148 & 0.347 & $+19.8$ pp \\
One child & 0.027 & 0.075 & $+4.7$ pp \\
Two children & 0.407 & 0.425 & $+1.8$ pp \\
Three or more & 0.417 & 0.153 & $-26.4$ pp \\
\bottomrule
\end{tabular}
\end{table}

This comparison characterises how the political equilibrium shifts when the household
problem includes fertility, not a policy counterfactual. The NIMBY benchmark has no
fertility margin---it is a different model, not a world where housing policy ignores
fertility. The within-model counterfactual is the scarcity experiment in
Section~\ref{sub:scarcity}, which compares the same fertility-aware economy under different
supply regimes.

The mechanism operates through both the direct and indirect housing channels identified in
the model section. The price-sensitive birth cost $\kappa_q q_t$ makes births more
expensive in high-price environments. The crowding margin $\lambda_c$ raises the shadow
cost of space when children are present. Together, these channels push first births later
and reduce the share of households who have children at all.

Figure~\ref{fig:first_birth_timing} plots the three timing objects against the full range
of steady-state house prices. The gradients are steep: moving from $q = 1.50$ to
$q = 3.00$ raises mean first-birth age by about 4.0~years and cuts the first-birth
rate by more than half. The wide-range comparative statics are reported in Appendix
Table~\ref{tab:timing_wide}.

\begin{figure}[htbp]
\centering
\includegraphics[width=\textwidth]{nimby_vs_fertility_first_birth_timing.png}
\caption{First-birth timing in the fertility benchmark. Higher house prices shift first births
to older ages, increase the share of first births occurring at age 30+, and lower the
first-birth rate.}
\label{fig:first_birth_timing}
\end{figure}

%% ============================================================
\subsection{Persistent housing scarcity and fertility}
\label{sub:scarcity}

A single political-tightness parameter is not enough for the claim that housing scarcity
lowers fertility, because it embeds the supply-tightening result in one dial that may be
doing too much work. A more disciplined comparison calibrates three different
supply-tightening shocks to the same housing outcome: a 15~percent higher house-price
index after 40~periods. The three margins are:
\begin{enumerate}
\item A reduction in supply elasticity ($\epsilon_0$ reduced by 6.2~percent);
\item A reduction in baseline supply ($s_0$ reduced by 2.4~percent);
\item A shift in baseline political tightness ($\theta_0$ increased by $+0.019$).
\end{enumerate}

This calibration strategy does two things. First, it shows that the result is not tied to
a single primitive. Second, it anchors the long-run exercise to a common housing target
rather than to an arbitrary parameter movement. Figure~\ref{fig:persistent_scarcity}
reports the resulting comparison.

Table~\ref{tab:scarcity_results} summarizes the quantitative outcomes. At horizon $t=40$,
all three shocks lower the fertility rate from 0.248 to 0.241---a reduction of roughly
0.007. The effect is modest but remarkably consistent across margins. On the historical
comparison, the calibrated shocks improve RMSE against the aggregate 1956--1995 fertility
decline by 3.4 to 4.0~percent and close roughly 5.4~percent of the baseline-to-data
end-gap.

\begin{table}[htbp]
\centering
\caption{Persistent housing scarcity: three supply-tightening margins}
\label{tab:scarcity_results}
\small
\begin{tabular}{l l cc c}
\toprule
Supply margin & Calibration & \multicolumn{2}{c}{Fertility at $t=40$} & RMSE \\
\cmidrule(lr){3-4}
 & & Baseline & Shock & improvement \\
\midrule
$\epsilon_0$ reduction & $\epsilon_0$ reduced 6.2\% & 0.248 & 0.241 & 4.0\% \\
$s_0$ reduction        & $s_0$ reduced 2.4\%        & 0.248 & 0.241 & 3.6\% \\
$\theta_0$ shift       & $\theta_0$ $+0.019$        & 0.248 & 0.241 & 3.4\% \\
\bottomrule
\end{tabular}
\end{table}

To put the 0.007 reduction in perspective: the model birth rate of 0.248 corresponds
roughly to a total fertility rate of 1.7 children per woman.%
\footnote{The transition birth rate of 0.248 is a population-weighted average across all age
groups, including those past childbearing. With the model's demographic structure, this maps
to approximately 1.7 births per woman over the reproductive window. The parity distribution
at the steady-state crossing implies completed fertility of approximately 1.46, which is
lower because the crossing is evaluated at a single price level rather than along the
transition path.} A 2.8~percent reduction
therefore translates to approximately 0.05 TFR points. Applied to the approximately
3.6~million annual US births, this implies roughly 100,000 fewer births per year
attributable to the housing-scarcity channel alone. Against the recent US fertility
decline (TFR fell from approximately 2.0 in 2007 to 1.62 in 2023, a drop of 0.38 points),
the housing channel accounts for roughly 13~percent of the observed decline---modest but
not negligible for a single mechanism operating through one margin of housing policy.

The effect also compounds over time. Table~\ref{tab:projection_results} reports the forward
projections under a medium-immigration demographic scenario. At the 40-period calibration
horizon, the fertility gap is 0.007 across all three margins. By 2050, the gap narrows
slightly as the direct shock effect stabilizes. But by 2100, the gap widens to
0.006--0.014, reflecting the interaction between housing scarcity and demographic aging:
restricted supply raises house prices further, which lowers realized fertility below what
aging alone would produce.

\begin{table}[htbp]
\centering
\caption{Compounding effect: fertility gap widens over time}
\label{tab:projection_results}
\small
\begin{tabular}{l cccc}
\toprule
 & Baseline & $\epsilon_0$ shock & $s_0$ shock & $\theta_0$ shock \\
\midrule
\multicolumn{5}{l}{\textit{Fertility rate}} \\
$t = 40$ (calibration) & 0.248 & 0.241 & 0.241 & 0.241 \\
2050 projection & 0.194 & 0.191 & 0.190 & 0.192 \\
2100 projection & 0.188 & 0.179 & 0.174 & 0.182 \\
\addlinespace
\multicolumn{5}{l}{\textit{House-price index}} \\
$t = 40$ (calibration) & 1.000 & 1.150 & 1.150 & 1.150 \\
2050 projection & 1.066 & 1.134 & 1.162 & 1.103 \\
2100 projection & 1.254 & 1.606 & 1.819 & 1.450 \\
\bottomrule
\end{tabular}
\end{table}

These projections are illustrative; at a 75-year horizon, structural changes in technology,
preferences, and policy make point forecasts unreliable.

The widening price gap in the bottom panel of Table~\ref{tab:projection_results} explains
the compounding: each supply margin that starts at a 15~percent price premium amplifies
over time, and the $s_0$ reduction reaches an 82~percent premium by 2100.

Housing scarcity therefore moves fertility in the right direction, and the effect grows
rather than dissipates. The model is better understood as identifying a compounding
channel---persistent underbuilding delays and discourages family formation, and the
demographic consequences feed back into housing demand---rather than as a complete account
of postwar fertility trends.

This also clarifies the relationship to the baby-boom exercise in Gross and Chivers
(2025). The baby-boom experiment asks how a larger cohort pushes on housing demand,
prices, and politics over time. The present paper asks the reverse question: how tighter
housing supply pushes back on fertility. The two mechanisms are therefore complementary,
not competing. A large cohort raises housing pressure through the children-at-home demand
channel; sustained housing scarcity then feeds back into later and lower realized
fertility. In that sense, the baby-boom transition is best read as a dynamic stress test
of the family-housing-demand mechanism, while the persistent-scarcity exercise is the main
causal direction emphasized here.

\begin{figure}[htbp]
\centering
\includegraphics[width=\textwidth]{nimby_shock_calibrated_supply_margins.png}
\caption{Persistent housing scarcity calibrated to a common housing target. Three different
supply-tightening shocks are chosen to raise the 40-period house-price index by 15 percent.
All three lower fertility.}
\label{fig:persistent_scarcity}
\end{figure}

%% ============================================================
\subsection{Voter preference comparison}
\label{sub:voter}

The fertility extension changes the political equilibrium because family-forming households
evaluate house-price increases differently from otherwise identical households without
children. Table~\ref{tab:voter_comparison} reports the net support index at selected price
levels in both the NIMBY benchmark (fertility shut off) and the fertility benchmark
(fertility active). The reported object is vote per unit mass, which ranges from $-1$ (all
households oppose higher prices) to $+1$ (all households support them). A positive value
means a net majority supports tighter supply restrictions at that price.

\begin{table}[htbp]
\centering
\caption{Net support for higher prices: NIMBY vs.\ fertility benchmarks}
\label{tab:voter_comparison}
\small
\begin{tabular}{l cc l}
\toprule
House price $q$ & NIMBY & Fertility & Interpretation \\
\midrule
1.50 & $+0.84$ & $+0.22$ & Both support higher prices; fertility support is much weaker \\
1.75 & $+0.79$ & $-0.00$ & NIMBY still supports; fertility is essentially at its local crossing \\
2.25 & $+0.23$ & $-0.28$ & NIMBY still supports; fertility opposes \\
2.50 & $-0.53$ & $-0.75$ & Both oppose; fertility more strongly \\
\bottomrule
\end{tabular}
\end{table}

The comparison reveals a sharp divergence, but the headline should come from the local
benchmark search rather than from a monotone reading of the coarse common grid. On the
common grid, households in the fertility benchmark are already close to indifferent at
$q = 1.75$ and oppose further price increases by $q = 2.25$, while the NIMBY benchmark
still supports higher prices at that point. The local search then places the political
crossing---the price at which the median voter switches from supporting to opposing tighter
supply---at 2.326 in the NIMBY benchmark versus 1.750 in the fertility benchmark.

The mechanism is straightforward. When household utility depends on family state, younger
households planning children place a higher value on affordable housing. The crowding
margin amplifies this: each additional child at home raises the shadow cost of housing
space, making price increases more costly. The fertility benchmark therefore shifts
political weight toward lower prices.

Two caveats apply. First, this is a within-model comparison, not an empirical claim about
actual voting behavior. The model does not include an explicit ballot or referendum
mechanism; it uses the same median-voter condition as Gross and Chivers (2025). Second, the
comparison assumes the same income process and demographic structure in both benchmarks. In
practice, the fertility decision itself could feed back into income and wealth accumulation
in ways the current model does not capture. The comparison is best read as showing that
endogenous fertility shifts the political equilibrium toward more permissive supply, not as
a prediction about the magnitude of that shift in any specific jurisdiction.

%% ============================================================
\subsection{Supporting evidence}
\label{sub:evidence}

The model implies that expensive housing should delay first births and lower fertility. The
paper does not rest on a settled empirical design, but two supporting pieces of evidence
point in the same direction as the model. The goal here is discipline and sign-checking, not
a definitive causal estimate.

Two pieces of evidence point in the same direction as the model. First, a panel of 28
European countries with joint OECD house-price and Eurostat first-birth data (486
country-year observations, country and year fixed effects) shows that higher real house
prices predict later first births: a one-unit increase in the OECD real house-price index
is associated with a 0.00033-year increase in the mean age at first birth ($p = 0.038$).
In a parallel specification covering 557 country-year observations, the same index predicts
lower total fertility ($\hat{\beta} = -0.00022$, $p = 0.006$). Both regressions include
country and year fixed effects, which absorb most of the overall variation. The implied
magnitudes are
small---a one-standard-deviation increase in the house-price index shifts mean first-birth
age by roughly 0.03~years---but the signs are correct and statistically significant. This
is best read as supportive evidence for the mechanism rather than as a tight calibration
target: within-country, within-year variation in house prices predicts the timing and level
of fertility in the direction the model implies, but the design is not the paper's main
contribution.

Second, the aggregate postwar comparison provides limited but informative context. Three
calibrated housing-scarcity shocks---each designed to raise the house-price index by
15~percent---improve the model's fit to the 1956--1995 fertility decline by 3.4 to
4.0~percent RMSE and close roughly 5.4~percent of the baseline-to-data end-gap.
Figure~\ref{fig:historical_validation} shows that the fertility extension fits the aligned
post-boom decline modestly better than the flat return path of the benchmark model. Housing
scarcity alone does not explain the full decline, but it adds a margin in the right
direction.

The overall evidence is consistent with the model but not yet conclusive. The cross-country
panel confirms the sign of the housing--fertility relationship but does not identify the
mechanism cleanly, since house prices covary with income, urbanisation, and educational
attainment. A fuller empirical design is left for future work. For the current paper, these
facts are best read as supporting context around a model-driven contribution rather than as a
fixed empirical core.

\begin{figure}[htbp]
\centering
\includegraphics[width=\textwidth]{nimby_vs_fertility_historical_validation.png}
\caption{Historical fertility validation. The fertility extension fits the aligned post-boom
decline modestly better than the flat return path of the benchmark model.}
\label{fig:historical_validation}
\end{figure}
